\documentclass[12pt]{article}
\usepackage[utf8]{inputenc}
\usepackage[T1]{fontenc}
\usepackage[margin=1in]{geometry}
\usepackage{setspace}
\usepackage{xcolor}
\usepackage[colorlinks=true, urlcolor=blue!60!black, linkcolor=blue!60!black]{hyperref}
\setlength{\parindent}{2em}
\setlength{\parskip}{1.25em}
\renewcommand{\baselinestretch}{1.0}
\begin{document}
\begin{center}
{\LARGE \textit{District of Columbia v. Heller}, 554 U.S. 570 (2008)\par}
\vspace{0.8em}
{\large Case Brief --- Four-Page Edition\par}
\end{center}
\vspace{1em}

\section{Memory Jogger}

In a 5--4 originalist decision, the Court held for the first time that the Second Amendment protects an individual right to possess and carry firearms unconnected with militia service, striking down the District of Columbia's handgun ban and trigger-lock mandate as infringements on self-defense in the home (\href{https://supreme.justia.com/cases/federal/us/554/570/}{554 U.S.\ 570}).

\section{Detailed Case Facts}

The District's Firearms Control Regulations Act of 1975 effectively banned civilian handgun possession by prohibiting handgun registration, barred carrying an unlicensed handgun even inside the home, and required lawful firearms such as rifles and shotguns to be kept ``unloaded and disassembled or bound by a trigger lock or similar device,'' with no exception for self-defense. Violations were crimes.

Dick Anthony Heller, a D.C.\ special police officer who carried a handgun on duty at the Federal Judicial Center, applied to register a handgun to keep at his home in a high-crime neighborhood, and the District refused. Because D.C.\ is a federal enclave, the case presented the Amendment's meaning without the incorporation question later resolved in \textit{McDonald v.\ City of Chicago} (2010). It came after nearly seventy years in which lower courts read \textit{United States v.\ Miller} (1939) as tying the right to militia service, a consensus broken by the Fifth Circuit's \textit{Emerson} decision and decades of scholarship and advocacy.

\section{Procedural History}

Six residents, including Heller, sued to enjoin the registration bar, the in-home licensing rule, and the trigger-lock mandate. The district court dismissed, finding no individual right outside militia service. The D.C.\ Circuit held that only Heller had standing and, 2--1, reversed on the merits. The Supreme Court granted certiorari, heard argument on March 18, 2008, and affirmed on June 26, 2008.

\section{Judicial Votes}

Affirmed, \textbf{5--4}.
\begin{itemize}
\item \textbf{Majority:} Justice Scalia, joined by Chief Justice Roberts and Justices Kennedy, Thomas, and Alito.
\item \textbf{Dissenting:} Justice Stevens, joined by Justices Souter, Ginsburg, and Breyer.
\item \textbf{Dissenting:} Justice Breyer, joined by Justices Stevens, Souter, and Ginsburg.
\end{itemize}

\section{Holding}

The Second Amendment guarantees a pre-existing individual right to keep and bear arms unconnected with militia service, with self-defense in the home at its core. Under any standard of scrutiny applicable to enumerated rights, a ban on an entire class of arms in common use for lawful self-defense and a mandate rendering firearms inoperable in the home are unconstitutional; the District must let Heller register and keep his handgun at home. The Court did not decide the validity of D.C.'s licensing regime or incorporation. Part III's ``safe harbor'' for felon and mental-illness bans, sensitive places, commercial conditions, and ``dangerous and unusual weapons'' is widely regarded as dicta, though later courts treat it as highly authoritative.

\section{Analysis of Opinions}

\textbf{Majority (Scalia).} Applying public-meaning originalism, Scalia divided the Amendment into a prefatory clause that announces a purpose and an operative clause that it cannot limit. ``The right of the people'' refers to individual rights, as in the First and Fourth Amendments; ``bear arms'' meant carrying weapons for confrontation, taking a military idiom only when followed by ``against.'' The militia meant all able-bodied men, and the clauses fit because tyrants destroyed militias by disarming the people, so the pre-existing right was codified to preserve the militia while self-defense remained its ``central component.'' State constitutions, early commentators, antebellum cases, and Reconstruction sources corroborated this reading. \textit{Cruikshank} and \textit{Presser} did not foreclose it, and \textit{Miller} protects only arms ``in common use at the time.'' The right is ``not unlimited,'' but rational-basis review and Breyer's interest balancing are rejected.

\textbf{Stevens (dissenting).} The Amendment protected the States' militias against federal disarmament, not private self-defense. Stevens faulted the majority for treating the preamble as surplusage, read ``keep and bear Arms'' as a unitary military phrase, stressed that Madison followed Virginia's military-focused proposal and rejected proposals protecting hunting or self-defense, and accused the Court of abandoning seventy years of reliance on \textit{Miller}.

\textbf{Breyer (dissenting).} Even assuming an individual self-defense interest, the ban was valid. Founding-era cities regulated guns, as in a 1783 Boston law barring loaded firearms in buildings. Rejecting strict scrutiny, Breyer proposed an ``interest-balancing inquiry'' and, given evidence of urban handgun homicides and legislative competence, upheld a ban that still allowed rifles and shotguns.

\section{Examples: Future Applications}

\begin{enumerate}
\item \textbf{Same-Side: Municipal Handgun Bans.} After incorporation, the Court struck down Chicago's ban in \textit{McDonald v.\ City of Chicago} (2010).
\item \textbf{Same-Side: ``Proper Cause'' Carry Licensing.} Extending \textit{Heller}'s text-and-history method beyond the home, \textit{Bruen} (2022) invalidated a ``special need'' requirement for public carry.
\item \textbf{Opposite-Side: Domestic-Violence Restraining Orders.} Disarming a person found to pose a credible threat to an intimate partner fits the tradition of disarming dangerous people and was upheld in \textit{United States v.\ Rahimi} (2024).
\item \textbf{Opposite-Side: Machine Guns and Short-Barreled Shotguns.} Under \textit{Heller}'s reading of \textit{Miller}, bans on ``dangerous and unusual weapons'' not in common use remain valid.
\item \textbf{Fence-Sitter: Assault-Weapon and Magazine Bans.} Challengers say millions own AR-15-style rifles, making them ``in common use''; defenders say their lethality makes them like M-16s and that handguns remain available.
\end{enumerate}

\section{Critique}

\textbf{Internal tensions and pragmatic objections.} Critics note that Scalia grounds the right in ``the people'' yet narrows it to ``law-abiding, responsible citizens,'' that the ``common use'' test is circular, and that the ``presumptively lawful'' list rests on no founding-era proof. Reva Siegel argues that \textit{Heller} reflects the modern gun-rights movement's popular constitutionalism, and Judges Posner and Wilkinson likened it to \textit{Roe}, faulting its history and its intrusion on local self-government.

\textbf{Originalist views.} Many originalists celebrate the decision as a victory for original public meaning over purposive balancing. Others, such as Nelson Lund, fault it for announcing no standard of review, which let lower courts apply intermediate scrutiny until \textit{Bruen}, and for unsupported exceptions in Part III.

\section{Key Quotations}

\begin{quote}
``Although we do not undertake an exhaustive historical analysis today of the full scope of the Second Amendment, nothing in our opinion should be taken to cast doubt on longstanding prohibitions on the possession of firearms by felons and the mentally ill, or laws forbidding the carrying of firearms in sensitive places such as schools and government buildings, or laws imposing conditions and qualifications on the commercial sale of arms.''
\end{quote}
\noindent --- Justice Scalia, Opinion of the Court, 554 U.S.\ at 626--627.

\begin{quote}
``Putting all of these textual elements together, we find that they guarantee the individual right to possess and carry weapons in case of confrontation.''
\end{quote}
\noindent --- Justice Scalia, Opinion of the Court, 554 U.S.\ at 592.

\begin{quote}
``We know of no other enumerated constitutional right whose core protection has been subjected to a freestanding `interest-balancing' approach. The very enumeration of the right takes out of the hands of government---even the Third Branch of Government---the power to decide on a case-by-case basis whether the right is really worth insisting upon. A constitutional guarantee subject to future judges' assessments of its usefulness is no constitutional guarantee at all.''
\end{quote}
\noindent --- Justice Scalia, Opinion of the Court, 554 U.S.\ at 634.

\end{document}
